% Generated draft: requires mathematical and source review.
\documentclass[11pt]{article}
\usepackage[margin=25mm]{geometry}
\usepackage{fontspec}
\setmainfont{Latin Modern Roman}
\usepackage{amsmath,amssymb,mathtools,unicode-math}
\setmathfont{Latin Modern Math}
\usepackage{enumitem,needspace,xurl,hyperref}
\setlength{\parindent}{0pt}
\setlength{\parskip}{5pt}
\setlength{\emergencystretch}{3em}
\newcommand{\N}{\mathbb{N}}
\newcommand{\Z}{\mathbb{Z}}
\newcommand{\Q}{\mathbb{Q}}
\newcommand{\R}{\mathbb{R}}
\newcommand{\C}{\mathbb{C}}
\newcommand{\F}{\mathbb{F}}
\newcommand{\A}{\mathbb{A}}
\newcommand{\PP}{\mathbb{P}}
\newcommand{\E}{\mathbb{E}}
\newcommand{\eps}{\varepsilon}
\newcommand{\dd}{\,\mathrm{d}}
\begin{document}
\section*{UnsolvedMath ID: 10400052}
\textbf{Source problem code:} catalogue entry
\subsection{Definitions and mathematical statement}
Let $\mathcal K^{\mathrm{or}}$ be the set of oriented smooth knots in the oriented three-sphere up to ambient isotopy respecting knot orientation. Use the Kontsevich invariant $Z$ with the source's unframed oriented-knot conventions, valued in the degree-completed diagram space. Equality means equality of every homogeneous-degree component, not equality after a fixed finite truncation.

The conjecture is injectivity:
\[
\forall K,K'\in\mathcal K^{\mathrm{or}},\qquad
Z(K)=Z(K')\quad\Longrightarrow\quad K=K'\text{ in }\mathcal K^{\mathrm{or}}.
\]
Orientation reversals or mirrors are not identified unless an actual oriented ambient isotopy realizes the equivalence.

In the characteristic-zero finite-type formulation supplied as equivalent, for every pair of distinct oriented knot types there should exist a finite-degree invariant $v$ with $v(K)\ne v(K')$. A knot invariant is of finite type of degree at most $d$ if its extension to singular knots by $v(K_\times)=v(K_+)-v(K_-)$ vanishes on every knot with $d+1$ transverse double points. The degree may depend on the pair; no uniform finite truncation separating every knot is asserted.

The equivalence uses the usual universality of the Kontsevich invariant for finite-type invariants, as referenced by the source's Conjecture 2.5. Precise diagram coefficient/normalization and framing conventions should match that source; different normalizations must not be compared as if equality were literal across definitions. Injectivity is not itself a finite stopping algorithm deciding knot equivalence.
\subsection{Short English statement}
Do all degrees of the Kontsevich invariant together distinguish oriented knot types? Equivalently, does some finite-type invariant separate every distinct pair, with its degree allowed to depend on the pair?
\subsection{Sources}
\begin{itemize}
\item Local UnsolvedMath registry, permanent ID 10400052, exact source code catalogue entry, ``Conjecture 3.2---The Kontsevich invariant distinguishes oriented knots''; no proposer/year supplied. Registry status open, classification OPEN-TRIAGE.
\item Ohtsuki, \emph{Problems on invariants of knots and 3-manifolds} (2002), Conjecture 3.2, PDF page 53, with Conjecture 2.5 cited there for the equivalent finite-type separation statement, \url{https://msp.org/gtm/2002/04/gtm-2002-04-024s.pdf}. Supplied pdf-font-structure extraction accessed July 29, 2026; provenance distinguishes 231 bold numbered items. No PDF was accessed here.
\item The supplied summary describes finite-type separation as a major open problem, without a later primary reference. This is supplied status assessment, not an independent current audit.
\item Provenance flags current primary-literature status and redistribution terms NEEDS\_REVIEW.
\end{itemize}
\end{document}
